% Options for packages loaded elsewhere
\PassOptionsToPackage{unicode}{hyperref}
\PassOptionsToPackage{hyphens}{url}
\PassOptionsToPackage{dvipsnames,svgnames,x11names}{xcolor}
%
\documentclass[
  11pt,
]{article}

\usepackage{amsmath,amssymb}
\usepackage{iftex}
\ifPDFTeX
  \usepackage[T1]{fontenc}
  \usepackage[utf8]{inputenc}
  \usepackage{textcomp} % provide euro and other symbols
\else % if luatex or xetex
  \usepackage{unicode-math}
  \defaultfontfeatures{Scale=MatchLowercase}
  \defaultfontfeatures[\rmfamily]{Ligatures=TeX,Scale=1}
\fi
\usepackage{lmodern}
\ifPDFTeX\else  
    % xetex/luatex font selection
    \setmainfont[]{Liberation Serif}
    \setsansfont[]{Liberation Sans}
\fi
% Use upquote if available, for straight quotes in verbatim environments
\IfFileExists{upquote.sty}{\usepackage{upquote}}{}
\IfFileExists{microtype.sty}{% use microtype if available
  \usepackage[]{microtype}
  \UseMicrotypeSet[protrusion]{basicmath} % disable protrusion for tt fonts
}{}
\makeatletter
\@ifundefined{KOMAClassName}{% if non-KOMA class
  \IfFileExists{parskip.sty}{%
    \usepackage{parskip}
  }{% else
    \setlength{\parindent}{0pt}
    \setlength{\parskip}{6pt plus 2pt minus 1pt}}
}{% if KOMA class
  \KOMAoptions{parskip=half}}
\makeatother
\usepackage{xcolor}
\usepackage[margin=1in]{geometry}
\setlength{\emergencystretch}{3em} % prevent overfull lines
\setcounter{secnumdepth}{-\maxdimen} % remove section numbering
% Make \paragraph and \subparagraph free-standing
\makeatletter
\ifx\paragraph\undefined\else
  \let\oldparagraph\paragraph
  \renewcommand{\paragraph}{
    \@ifstar
      \xxxParagraphStar
      \xxxParagraphNoStar
  }
  \newcommand{\xxxParagraphStar}[1]{\oldparagraph*{#1}\mbox{}}
  \newcommand{\xxxParagraphNoStar}[1]{\oldparagraph{#1}\mbox{}}
\fi
\ifx\subparagraph\undefined\else
  \let\oldsubparagraph\subparagraph
  \renewcommand{\subparagraph}{
    \@ifstar
      \xxxSubParagraphStar
      \xxxSubParagraphNoStar
  }
  \newcommand{\xxxSubParagraphStar}[1]{\oldsubparagraph*{#1}\mbox{}}
  \newcommand{\xxxSubParagraphNoStar}[1]{\oldsubparagraph{#1}\mbox{}}
\fi
\makeatother


\providecommand{\tightlist}{%
  \setlength{\itemsep}{0pt}\setlength{\parskip}{0pt}}\usepackage{longtable,booktabs,array}
\usepackage{calc} % for calculating minipage widths
% Correct order of tables after \paragraph or \subparagraph
\usepackage{etoolbox}
\makeatletter
\patchcmd\longtable{\par}{\if@noskipsec\mbox{}\fi\par}{}{}
\makeatother
% Allow footnotes in longtable head/foot
\IfFileExists{footnotehyper.sty}{\usepackage{footnotehyper}}{\usepackage{footnote}}
\makesavenoteenv{longtable}
\usepackage{graphicx}
\makeatletter
\newsavebox\pandoc@box
\newcommand*\pandocbounded[1]{% scales image to fit in text height/width
  \sbox\pandoc@box{#1}%
  \Gscale@div\@tempa{\textheight}{\dimexpr\ht\pandoc@box+\dp\pandoc@box\relax}%
  \Gscale@div\@tempb{\linewidth}{\wd\pandoc@box}%
  \ifdim\@tempb\p@<\@tempa\p@\let\@tempa\@tempb\fi% select the smaller of both
  \ifdim\@tempa\p@<\p@\scalebox{\@tempa}{\usebox\pandoc@box}%
  \else\usebox{\pandoc@box}%
  \fi%
}
% Set default figure placement to htbp
\def\fps@figure{htbp}
\makeatother

\makeatletter
\@ifpackageloaded{tcolorbox}{}{\usepackage[skins,breakable]{tcolorbox}}
\@ifpackageloaded{fontawesome5}{}{\usepackage{fontawesome5}}
\definecolor{quarto-callout-color}{HTML}{909090}
\definecolor{quarto-callout-note-color}{HTML}{0758E5}
\definecolor{quarto-callout-important-color}{HTML}{CC1914}
\definecolor{quarto-callout-warning-color}{HTML}{EB9113}
\definecolor{quarto-callout-tip-color}{HTML}{00A047}
\definecolor{quarto-callout-caution-color}{HTML}{FC5300}
\definecolor{quarto-callout-color-frame}{HTML}{acacac}
\definecolor{quarto-callout-note-color-frame}{HTML}{4582ec}
\definecolor{quarto-callout-important-color-frame}{HTML}{d9534f}
\definecolor{quarto-callout-warning-color-frame}{HTML}{f0ad4e}
\definecolor{quarto-callout-tip-color-frame}{HTML}{02b875}
\definecolor{quarto-callout-caution-color-frame}{HTML}{fd7e14}
\makeatother
\makeatletter
\@ifpackageloaded{caption}{}{\usepackage{caption}}
\AtBeginDocument{%
\ifdefined\contentsname
  \renewcommand*\contentsname{Table of contents}
\else
  \newcommand\contentsname{Table of contents}
\fi
\ifdefined\listfigurename
  \renewcommand*\listfigurename{List of Figures}
\else
  \newcommand\listfigurename{List of Figures}
\fi
\ifdefined\listtablename
  \renewcommand*\listtablename{List of Tables}
\else
  \newcommand\listtablename{List of Tables}
\fi
\ifdefined\figurename
  \renewcommand*\figurename{Figure}
\else
  \newcommand\figurename{Figure}
\fi
\ifdefined\tablename
  \renewcommand*\tablename{Table}
\else
  \newcommand\tablename{Table}
\fi
}
\@ifpackageloaded{float}{}{\usepackage{float}}
\floatstyle{ruled}
\@ifundefined{c@chapter}{\newfloat{codelisting}{h}{lop}}{\newfloat{codelisting}{h}{lop}[chapter]}
\floatname{codelisting}{Listing}
\newcommand*\listoflistings{\listof{codelisting}{List of Listings}}
\makeatother
\makeatletter
\makeatother
\makeatletter
\@ifpackageloaded{caption}{}{\usepackage{caption}}
\@ifpackageloaded{subcaption}{}{\usepackage{subcaption}}
\makeatother

\usepackage{bookmark}

\IfFileExists{xurl.sty}{\usepackage{xurl}}{} % add URL line breaks if available
\urlstyle{same} % disable monospaced font for URLs
\hypersetup{
  pdftitle={Lab 2: Frequency Distributions and Graphs},
  colorlinks=true,
  linkcolor={blue},
  filecolor={Maroon},
  citecolor={Blue},
  urlcolor={Blue},
  pdfcreator={LaTeX via pandoc}}


\title{Lab 2: Frequency Distributions and Graphs}
\usepackage{etoolbox}
\makeatletter
\providecommand{\subtitle}[1]{% add subtitle to \maketitle
  \apptocmd{\@title}{\par {\large #1 \par}}{}{}
}
\makeatother
\subtitle{Covers Chapter 2: Frequency Distributions \textbar{} 50-minute
lab \textbar{} Continues from Lab 1}
\author{}
\date{}

\begin{document}
\maketitle


\textbf{Before you start:} Open SPSS, then open your saved file from
last week: \textbf{File → Open → Data} → navigate to where you saved it
→ select \texttt{Lab1\_YourLastName.sav} → \textbf{Open}. You should see
your 10-row, 6-variable dataset (StudentID, Major, ClassYear,
StudyHours, NetflixHrs, Stress). We are \emph{not} creating new
variables today --- we're analyzing the ones you already built.

\begin{center}\rule{0.5\linewidth}{0.5pt}\end{center}

\subsection{Part 1: Frequency Table + Bar Chart for a Categorical
Variable (13
min)}\label{part-1-frequency-table-bar-chart-for-a-categorical-variable-13-min}

We'll start with \textbf{Major} --- a nominal variable --- and build
both a frequency table and its graph in one pass.

\begin{enumerate}
\def\labelenumi{\arabic{enumi}.}
\tightlist
\item
  \textbf{Analyze → Descriptive Statistics → Frequencies}
\item
  Click \textbf{Major} in the left-hand list, click the arrow to move it
  into the \textbf{Variable(s)} box.
\item
  Click the \textbf{Charts\ldots{}} button.

  \begin{itemize}
  \tightlist
  \item
    Select the \textbf{Bar charts} radio button.
  \item
    Under ``Chart Values,'' select \textbf{Percentages} (instead of the
    default ``Frequencies'') --- this shows the proportion of students
    in each major, not just the raw count.
  \item
    Click \textbf{Continue}.
  \end{itemize}
\item
  Click \textbf{OK}.
\end{enumerate}

Two things appear in the Output Viewer:

\textbf{A frequency table.} Match its columns to the textbook's
vocabulary:

\begin{longtable}[]{@{}
  >{\raggedright\arraybackslash}p{(\linewidth - 2\tabcolsep) * \real{0.5000}}
  >{\raggedright\arraybackslash}p{(\linewidth - 2\tabcolsep) * \real{0.5000}}@{}}
\toprule\noalign{}
\begin{minipage}[b]{\linewidth}\raggedright
SPSS column
\end{minipage} & \begin{minipage}[b]{\linewidth}\raggedright
Chapter 2 term
\end{minipage} \\
\midrule\noalign{}
\endhead
\bottomrule\noalign{}
\endlastfoot
Frequency & \emph{f} --- the raw count in each category \\
Percent & \% --- same as computing p × 100, where p = f/N \\
Valid Percent & \% excluding any missing data (identical to Percent
here, since we have none) \\
Cumulative Percent & running total as you move down the table --- this
is the basis for percentile ranks \\
\end{longtable}

\textbf{A bar chart.} Notice the spaces between the bars --- that's
deliberate. Chapter 2 draws bar charts with gaps specifically for
nominal/ordinal data, to visually signal that the categories are
separate and unordered, unlike a histogram's touching bars for numerical
data.

\begin{tcolorbox}[enhanced jigsaw, arc=.35mm, breakable, opacityback=0, colframe=quarto-callout-tip-color-frame, left=2mm, colbacktitle=quarto-callout-tip-color!10!white, opacitybacktitle=0.6, bottomrule=.15mm, toprule=.15mm, rightrule=.15mm, coltitle=black, bottomtitle=1mm, titlerule=0mm, toptitle=1mm, leftrule=.75mm, title=\textcolor{quarto-callout-tip-color}{\faLightbulb}\hspace{0.5em}{Checkpoint}, colback=white]

In your output, which major has the highest frequency? What percentage
of the class does it represent?

\end{tcolorbox}

\begin{center}\rule{0.5\linewidth}{0.5pt}\end{center}

\subsection{Part 2: Independent Practice --- ClassYear (5
min)}\label{part-2-independent-practice-classyear-5-min}

Repeat the same process for \textbf{ClassYear}, an ordinal variable, on
your own:

\begin{enumerate}
\def\labelenumi{\arabic{enumi}.}
\tightlist
\item
  \textbf{Analyze → Descriptive Statistics → Frequencies} → move
  \textbf{ClassYear} in → \textbf{Charts} → \textbf{Bar charts} →
  \textbf{Percentages} → \textbf{Continue} → \textbf{OK}.
\end{enumerate}

\begin{tcolorbox}[enhanced jigsaw, arc=.35mm, breakable, opacityback=0, colframe=quarto-callout-tip-color-frame, left=2mm, colbacktitle=quarto-callout-tip-color!10!white, opacitybacktitle=0.6, bottomrule=.15mm, toprule=.15mm, rightrule=.15mm, coltitle=black, bottomtitle=1mm, titlerule=0mm, toptitle=1mm, leftrule=.75mm, title=\textcolor{quarto-callout-tip-color}{\faLightbulb}\hspace{0.5em}{Checkpoint}, colback=white]

Looking at the Cumulative Percent column, what percentage of students
are Sophomores or below? (This is exactly what a percentile rank tells
you --- the percentage of the group at or below a given point.)

\end{tcolorbox}

\begin{center}\rule{0.5\linewidth}{0.5pt}\end{center}

\subsection{Part 3: A Numerical Variable --- and Why It Gets Messy (8
min)}\label{part-3-a-numerical-variable-and-why-it-gets-messy-8-min}

Now let's try the same recipe on a \textbf{scale (numerical)} variable:
\textbf{StudyHours}.

\begin{enumerate}
\def\labelenumi{\arabic{enumi}.}
\tightlist
\item
  \textbf{Analyze → Descriptive Statistics → Frequencies}
\item
  Remove Major/ClassYear if they're still listed; move
  \textbf{StudyHours} in instead.
\item
  Click \textbf{Statistics\ldots{}} and check the box for
  \textbf{Quartiles}. Click \textbf{Continue}.

  \begin{itemize}
  \tightlist
  \item
    Quartiles report the 25th, 50th, and 75th percentiles --- this is
    SPSS computing exactly the percentile concept from this chapter.
  \end{itemize}
\item
  Click \textbf{OK} (leave Charts unchecked for now).
\end{enumerate}

\textbf{Look at the frequency table.} With only a few exceptions, every
row has a frequency of 1 --- because with continuous data spread across
a wide range, individual scores rarely repeat.

This is precisely the problem Chapter 2 describes as the reason for
\textbf{grouped frequency distribution tables}: when scores cover a wide
range, listing every individual value produces a long, uninformative
table instead of a picture you can actually interpret.

\begin{tcolorbox}[enhanced jigsaw, arc=.35mm, breakable, opacityback=0, colframe=quarto-callout-tip-color-frame, left=2mm, colbacktitle=quarto-callout-tip-color!10!white, opacitybacktitle=0.6, bottomrule=.15mm, toprule=.15mm, rightrule=.15mm, coltitle=black, bottomtitle=1mm, titlerule=0mm, toptitle=1mm, leftrule=.75mm, title=\textcolor{quarto-callout-tip-color}{\faLightbulb}\hspace{0.5em}{Checkpoint}, colback=white]

Look at the Percentile Values table in your output. What StudyHours
value corresponds to the 50th percentile?

\end{tcolorbox}

\begin{center}\rule{0.5\linewidth}{0.5pt}\end{center}

\subsection{Part 4: Building a Grouped Frequency Table (12
min)}\label{part-4-building-a-grouped-frequency-table-12-min}

We'll now group StudyHours into class intervals, following the
textbook's guidelines: about 3--10 intervals (we have a small dataset,
so we'll aim for 3), a simple interval width (5 is a good round number),
and a bottom score that's a multiple of the width.

\textbf{Why the boundaries below look slightly unusual:} we're setting
the cut points at the \emph{real limits} (4.5, 9.5, 14.5, 19.5) rather
than the whole numbers themselves. This matters because our data
includes decimals (like 9.5 hours) --- using the real limit means every
possible score falls into exactly one interval with no gaps, exactly as
Chapter 2 describes.

\begin{enumerate}
\def\labelenumi{\arabic{enumi}.}
\tightlist
\item
  \textbf{Transform → Recode into Different Variables}
\item
  Click \textbf{StudyHours}, move it into the middle box (it will show
  as \texttt{StudyHours\ -\/-\textgreater{}\ ?}).
\item
  In the \textbf{Output Variable} box on the right: type
  \texttt{StudyGroup} in \textbf{Name}, type
  \texttt{Study\ hours\ (grouped)} in \textbf{Label}. Click
  \textbf{Change}.
\item
  Click \textbf{Old and New Values\ldots{}} A new dialog opens. Set up
  three ranges, clicking \textbf{Add} after each:

  \begin{itemize}
  \tightlist
  \item
    Old Value → select \textbf{Range}: \texttt{4.5} through
    \texttt{9.4}. New Value: \texttt{1}. Click \textbf{Add}.
  \item
    Old Value → select \textbf{Range}: \texttt{9.5} through
    \texttt{14.4}. New Value: \texttt{2}. Click \textbf{Add}.
  \item
    Old Value → select \textbf{Range}: \texttt{14.5} through
    \texttt{19.4}. New Value: \texttt{3}. Click \textbf{Add}.
  \end{itemize}
\item
  Click \textbf{Continue}, then \textbf{OK}. A new column,
  \texttt{StudyGroup}, appears at the end of your Data View.
\item
  Switch to \textbf{Variable View}. Find the \texttt{StudyGroup} row and
  finish setting it up like you did in Lab 1:

  \begin{itemize}
  \tightlist
  \item
    \textbf{Decimals:} \texttt{0}
  \item
    \textbf{Values:} \texttt{1} = \texttt{5-9\ hrs}, \texttt{2} =
    \texttt{10-14\ hrs}, \texttt{3} = \texttt{15-19\ hrs}
  \item
    \textbf{Measure:} select \textbf{Ordinal} (the groups now have a
    meaningful order, even though the underlying variable is scale)
  \end{itemize}
\item
  Run \textbf{Analyze → Descriptive Statistics → Frequencies} on
  \textbf{StudyGroup}.
\end{enumerate}

\textbf{Compare this table to Part 3's.} Same underlying data, but now
it tells a story you can actually see at a glance --- that's the entire
purpose of grouping.

\begin{center}\rule{0.5\linewidth}{0.5pt}\end{center}

\subsection{Part 5: Histogram and Shape (6
min)}\label{part-5-histogram-and-shape-6-min}

\begin{enumerate}
\def\labelenumi{\arabic{enumi}.}
\tightlist
\item
  \textbf{Analyze → Descriptive Statistics → Frequencies}
\item
  Move \textbf{StudyHours} (the original, ungrouped variable) into the
  Variable(s) box.
\item
  Click \textbf{Charts\ldots{}} → select \textbf{Histograms} → click
  \textbf{Continue} → \textbf{OK}.
\end{enumerate}

SPSS automatically bins the continuous variable for you (similar in
spirit to what you just did by hand in Part 4). Look at the shape of the
bars.

\begin{tcolorbox}[enhanced jigsaw, arc=.35mm, breakable, opacityback=0, colframe=quarto-callout-tip-color-frame, left=2mm, colbacktitle=quarto-callout-tip-color!10!white, opacitybacktitle=0.6, bottomrule=.15mm, toprule=.15mm, rightrule=.15mm, coltitle=black, bottomtitle=1mm, titlerule=0mm, toptitle=1mm, leftrule=.75mm, title=\textcolor{quarto-callout-tip-color}{\faLightbulb}\hspace{0.5em}{Checkpoint --- answer using your histogram}, colback=white]

\begin{itemize}
\tightlist
\item
  Is the distribution roughly \textbf{symmetrical}, or does it look
  \textbf{skewed}?
\item
  If skewed, which direction is the tail pointing --- toward the higher
  scores (positive skew) or the lower scores (negative skew)?
\item
  Take a look at where SPSS started the Y-axis. Does it start at zero?
  (Keep this in mind --- we'll come back to it below.)
\end{itemize}

\end{tcolorbox}

\begin{center}\rule{0.5\linewidth}{0.5pt}\end{center}

\subsection{Take-Home Lab Report (due before next
class)}\label{take-home-lab-report-due-before-next-class}

\begin{enumerate}
\def\labelenumi{\arabic{enumi}.}
\tightlist
\item
  \textbf{By-hand check.} Using the Major frequency table from Part 1,
  compute Σ\emph{fX} by hand --- but first you'll need to assign a
  numeric code to each major (use the same codes from your Values
  labels: 1=Psychology, 2=Biology, 3=Business, 4=Undecided) and treat
  frequency × code as \emph{fX} for each row. Show your work. (This is
  an artificial calculation since Major is nominal and its codes aren't
  truly quantitative --- that's the point: explain in one sentence why
  Σ\emph{fX} is mathematically possible here but not meaningful.)
\item
  \textbf{Polygon by hand.} Using your Part 4 grouped frequency table
  (StudyGroup), sketch a frequency polygon on paper: place a dot above
  the \emph{midpoint} of each interval at a height equal to its
  frequency, connect the dots, and drop the line to zero one
  interval-width beyond each end. (Recall: the midpoint of a class
  interval is the average of its highest and lowest apparent limits.)
\item
  \textbf{Stem-and-leaf by hand.} Construct a stem-and-leaf display for
  the 10 StudyHours values. What advantage does this display have over
  the grouped frequency table you built in Part 4?
\item
  \textbf{Population vs.~sample.} Our 10 students are a sample. If we
  had data on the \emph{entire population} of students at this
  university instead, would we typically graph actual frequencies or
  relative frequencies? Why does the textbook recommend relative
  frequencies for very large groups?
\item
  \textbf{The Y-axis question.} Referring to Box 2.1 in the textbook
  (the homicide-rate graph example) and what you noticed at the end of
  Part 5: why does it matter whether a graph's Y-axis starts at zero?
  Give a one-sentence explanation of how a manipulated Y-axis could
  mislead someone looking at your StudyHours histogram.
\item
  \textbf{Percentile interpretation.} In your own words, explain what it
  means that a StudyHours score has a percentile rank of 70\%. Use the
  cumulative percent column from one of today's frequency tables as your
  example.
\end{enumerate}

\begin{center}\rule{0.5\linewidth}{0.5pt}\end{center}

\subsection{Instructor Notes}\label{instructor-notes}

\begin{tcolorbox}[enhanced jigsaw, arc=.35mm, breakable, opacityback=0, colframe=quarto-callout-important-color-frame, left=2mm, colbacktitle=quarto-callout-important-color!10!white, opacitybacktitle=0.6, bottomrule=.15mm, toprule=.15mm, rightrule=.15mm, coltitle=black, bottomtitle=1mm, titlerule=0mm, toptitle=1mm, leftrule=.75mm, title=\textcolor{quarto-callout-important-color}{\faExclamation}\hspace{0.5em}{Trim this section before distributing to students}, colback=white]

\begin{itemize}
\tightlist
\item
  \textbf{This is the tightest lab yet.} Part 4 (Recode) is the most
  likely place to run over --- if the section is behind schedule,
  consider demonstrating Part 4 live on the projector rather than having
  students execute it themselves, and shift the hands-on grouping
  practice to the take-home assignment using a walkthrough video
  instead.
\item
  \textbf{The 4.5/9.5/14.5 boundary values are deliberate}, not
  arbitrary --- they're the real limits, and they're chosen so that no
  student's actual StudyHours value (6.0, 7.5, 8.5, 9.5, 10.0, 11.0,
  12.5, 13.0, 15.0, 18.5) sits ambiguously on a boundary. If you swap in
  a different class dataset, re-derive the cut points the same way
  rather than reusing 4.5/9.5/14.5 by default.
\item
  \textbf{Common error:} students select ``Value'' instead of ``Range''
  in the Recode dialog, or forget to click ``Add'' after each range
  before moving to the next --- same failure mode as the Values dialog
  in Lab 1.
\item
  \textbf{Shape discussion (Part 5):} with only 10 data points, resist
  letting students over-claim a strong skew --- a good moment to
  explicitly note that shape judgments are much shakier with small
  samples, which sets up why later chapters (central tendency,
  variability) give us more precise numerical tools than ``eyeballing''
  a graph.
\item
  Take-home Question 1 is intentionally a little bit of a trap ---
  Σ\emph{fX} is computable for nominal data but meaningless, and
  catching that is the actual learning objective, not the arithmetic.
\end{itemize}

\end{tcolorbox}




\end{document}
